\section{Morphismes équidimensionnels}
\label{section:IV.13-fr}

Ce paragraphe est consacré à l'étude de la variation de la \emph{dimension}
des fibres d'un morphisme localement de type fini $f:X\to Y$ (qui est déjà
intervenue à propos de la «~formule des dimensions~» dans 5.5 et 5.6). On
prouve d'abord (13.1.3) que la fonction
$x\mapsto\dim_x(f^{-1}(f(x)))$ est toujours \emph{semi-continue
supérieurement} dans $X$ (théorème de semi-continuité de Chevalley). On étudie
ensuite plus particulièrement les morphismes, dits «~équidimensionnels~», pour
lesquels cette fonction est \emph{localement constante}.

\oldpage[IV-3]{188}
Malheureusement la notion de morphisme équidimensionnel n'est pas stable par
changement de base~; c'est pourquoi dans de nombreuses questions il est plus
commode de travailler avec la notion de morphisme universellement ouvert, dont
l'étude fait l'objet des \S\S~14 et 15.

\subsection{Le théorème de semi-continuité de Chevalley}
\label{subsection:IV.13.1-fr}

\begin{lemma}[13.1.1]
\label{IV.13.1.1-fr}
Soient $Y$ un préschéma irréductible localement noethérien, $X$ un préschéma
irréductible, $f:X\to Y$ un morphisme dominant de type fini. Soit $\xi$
(resp. $\eta$) le point générique de $X$ (resp. $Y$) et soit
$e=\dim(f^{-1}(\eta))$. Alors, pour tout $x\in X$, toutes les composantes
irréductibles de $f^{-1}(f(x))$ sont de dimension $\geq e$.
\end{lemma}

\begin{proof}
La proposition est immédiate lorsque, pour tout $y\in f(X)$,
$\mathcal O_y$ est un \emph{anneau universellement caténaire}~: en effet, si
$x$ est point générique d'une composante irréductible $Z$ de $f^{-1}(y)$, il
résulte de (5.6.5), joint à (5.2.1), que l'on a
$e+\dim(\mathcal O_y)=\dim(Z)+\dim(\mathcal O_x)$~; mais en vertu de
($\mathbf 0$, 16.3.9), on a $\dim(\mathcal O_x)\leq\dim(\mathcal O_y)$,
d'où la conclusion dans ce cas.

Nous allons ramener le cas général à ce cas particulier. La question est
évidemment locale sur $Y$, et, compte tenu de (4.1.1.3), elle est aussi locale
sur $X$~; on peut donc se borner au cas où $Y=\operatorname{Spec}(A)$ et
$X=\operatorname{Spec}(B)$ sont affines et irréductibles, $B$ étant une
$A$-algèbre de type fini. En outre (1.5.4), on peut supposer $X$ et $Y$
réduits, donc $A$ et $B$ \emph{intègres} et, comme $f$ est dominant, $A$ est
alors un \emph{sous-anneau} de $B$. Considérons $A$ comme limite inductive de
ses sous-$\mathbf Z$-algèbres de type fini~; il résulte alors de (8.9.1) qu'il
existe une telle sous-algèbre $A_0$ et une $A_0$-algèbre de type fini $B_0$
telles que $B=B_0\otimes_{A_0}A$. Posons $Y_0=\operatorname{Spec}(A_0)$,
$X_0=\operatorname{Spec}(B_0)$, et soit $f_0:X_0\to Y_0$ le morphisme
correspondant à l'homomorphisme $A_0\to B_0$, de sorte que
$f=f_0\times 1_Y$. Il n'est pas évident \emph{a priori} que le préschéma $X_0$
soit intègre, mais nous allons voir qu'on peut se ramener à ce cas. Soit
$\eta_0$ le point générique de $Y_0$, de sorte que si $g$ est le morphisme
$Y\to Y_0$, on a $g(\eta)=\eta_0$~; par transitivité des fibres
(\textbf{I}, 3.6.4), on a
$f^{-1}(\eta)=f_0^{-1}(\eta_0)\otimes_{k(\eta_0)}k(\eta)$, et comme
$f^{-1}(\eta)$ est irréductible par hypothèse, il en est de même de
$f_0^{-1}(\eta_0)$ (4.4.1). Notre assertion résultera alors du lemme suivant~:

\begin{lemma}[13.1.2]
\label{IV.13.1.2-fr}
Soient $Y_0$, $Y$ deux préschémas intègres de points génériques $\eta_0$,
$\eta$, $g:Y\to Y_0$ un morphisme dominant, $f_0:X_0\to Y_0$ un morphisme
dominant tel que $f_0^{-1}(\eta_0)$ soit irréductible. Soit $X'_0$ l'unique
composante irréductible de $X_0$ rencontrant $f_0^{-1}(\eta_0)$
($\mathbf 0_{\mathrm I}$, 2.1.8), et désignons encore par $X'_0$ le
sous-préschéma fermé réduit de $X_0$ ayant $X'_0$ pour espace sous-jacent.
Supposons que le préschéma $X=X_0\times_{Y_0}Y$ soit intègre~; alors $X$ est
isomorphe à $X'_0\times_{Y_0}Y$.
\end{lemma}

En effet, si $j_0:X'_0\to X_0$ est l'injection canonique, qui est une immersion
fermée, $j=j_0\times 1_Y:X'_0\times_{Y_0}Y\to X_0\times_{Y_0}Y=X$ est une
immersion fermée. D'autre part, $X'_0$ contient la fibre
$f_0^{-1}(\eta_0)$, donc $X'_0\times_{Y_0}Y$ contient $f^{-1}(\eta)$~;
notons d'autre part que $f^{-1}(\eta)$ n'est pas vide (\textbf{I}, 3.4.7),
donc contient le point générique de $X$~; par suite l'image de $j$ est
nécessairement $X$ tout entier. Mais comme $X$ est intègre, le seul
sous-préschéma fermé de $X$ ayant $X$ pour espace sous-jacent est $X$
lui-même, donc $j$ est un isomorphisme.

Ce lemme étant établi, on peut donc supposer que $X_0$ est intègre~; pour tout
$y_0\in Y_0$, $\mathcal O_{y_0}$ est une $\mathbf Z$-algèbre essentiellement
de type fini, donc un anneau universellement

\oldpage[IV-3]{189}
caténaire (5.6.4)~; par suite, pour tout $y_0\in f_0(X_0)$, toutes les
composantes irréductibles de $f_0^{-1}(y)$ ont une dimension $\geq e$, car on
sait que $\dim(f_0^{-1}(\eta_0))=e$ par transitivité des fibres (4.1.4). Pour
tout $y\in f(X)$, on a alors $y_0=g(y)\in f_0(X_0)$ et par transitivité des
fibres et (4.2.7), on achève la démonstration.
\end{proof}

\begin{theorem}[13.1.3]
\label{IV.13.1.3-fr}
\textnormal{(Chevalley).} --- Soit $f:X\to Y$ un morphisme localement de type
fini. Pour tout entier $n$, l'ensemble $F_n(X)$ des $x\in X$ tels que
$\dim_x(f^{-1}(f(x)))\geq n$ est fermé~; autrement dit, la fonction
$x\mapsto\dim_x(f^{-1}(f(x)))$ est semi-continue supérieurement dans $X$.
\end{theorem}

\begin{proof}
I) Supposons d'abord que $f$ soit localement de présentation finie.

La question est évidemment locale sur $X$ et sur $Y$, et on peut donc supposer
$Y=\operatorname{Spec}(A)$, $X=\operatorname{Spec}(B)$ affines, $B$ étant une
$A$-algèbre de présentation finie. On sait alors (8.9.1) qu'il y a un
sous-anneau noethérien $A_0$ de $A$, et une $A_0$-algèbre de type fini $B_0$
tels que si l'on pose $Y_0=\operatorname{Spec}(A_0)$,
$X_0=\operatorname{Spec}(B_0)$, on ait

\[
  X=X_0\times_{Y_0}Y,\qquad f=f_0\times 1_Y,
\]

où $f_0:X_0\to Y_0$ correspond à l'homomorphisme $A_0\to B_0$. Soient
$g:Y\to Y_0$ le morphisme correspondant à l'injection canonique $A_0\to A$,
$y$ un point de $Y$, $y_0=g(y)$~; on sait que
$f^{-1}(y)=f_0^{-1}(y_0)\otimes_{k(y_0)}k(y)$, et il résulte de (4.2.7) que
si $p$ est la projection canonique $f^{-1}(y)\to f_0^{-1}(y_0)$, les
composantes irréductibles de $f^{-1}(y)$ sont les composantes irréductibles des
ensembles $p^{-1}(Z_0)$, où $Z_0$ parcourt l'ensemble des composantes
irréductibles de $f_0^{-1}(y_0)$, et chacune des composantes irréductibles de
$p^{-1}(Z_0)$ domine $Z_0$ et a une dimension égale à $\dim(Z_0)$. Tenant
compte de ($\mathbf 0$, 14.1.5), on voit donc que si $g':X\to X_0$ est la
projection canonique, $x$ un point de $X$ et $x_0=g'(x)$, on a
$\dim_x(f^{-1}(f(x)))=\dim_{x_0}(f_0^{-1}(f_0(x_0)))$, d'où
$F_n(X)=g'^{-1}(F_n(X_0))$, et on est par suite ramené à démontrer le théorème
lorsque $Y$ est \emph{noethérien}, ce que nous supposerons désormais.

On peut évidemment supposer $X$ et $Y$ réduits (1.5.4). En considérant
l'ensemble des parties fermées $Y'$ de $Y$ telles que le théorème soit vrai
pour le sous-préschéma fermé de $Y$ ayant $Y'$ pour espace sous-jacent et pour
$X'=f^{-1}(Y')$, on peut raisonner par récurrence noethérienne
($\mathbf 0_{\mathrm I}$, 2.2.2), et supposer que le théorème est vrai pour
toute partie fermée $Y'\neq Y$ de $Y$. Si $X_i$ ($1\leq i\leq m$) sont les
sous-préschémas fermés réduits de $X$ ayant pour espaces sous-jacents les
composantes irréductibles de $X$, on a $F_n(X)=\bigcup_iF_n(X_i)$ en vertu de
($\mathbf 0$, 14.1.5), et l'on peut donc se borner à démontrer le théorème pour
chacun des $X_i$, autrement dit, on peut supposer $X$ irréductible. Si $Z$ est
le sous-préschéma fermé de $Y$ ayant $\overline{f(X)}$ pour espace sous-jacent,
$f$ se factorise en $X\xrightarrow{g}Z\xrightarrow{j}Y$, où $j$ est
l'injection canonique (\textbf{I}, 5.2.2), et $g$ est de type fini (1.5.4),
donc il suffit de démontrer le théorème pour $Z$ et $g$~; en vertu de
l'hypothèse de récurrence, on est donc ramené à considérer seulement le cas où
$Z=Y$, autrement dit où $Y$ est irréductible et $f$ dominant. Soit alors
$\eta$ le point générique de $Y$ et posons $e=\dim(f^{-1}(\eta))$~; il résulte
de (13.1.1) que pour $n\leq e$ on a $F_n(X)=X$, et par suite on peut se borner
au cas où $n>e$. Mais alors (9.5.6), il y a un voisinage ouvert $U$ de $\eta$
dans $Y$ tel que $F_n(X)\subset f^{-1}(Y-U)$~; comme $Y-U\neq Y$, l'hypothèse
de récurrence entraîne que $F_n(X)$ est fermé.

II) Passons maintenant au cas général, en supposant toujours
$S=\operatorname{Spec}(A)$ et $X=\operatorname{Spec}(B)$ affines, $B$ étant
une $A$-algèbre de type fini, donc de la forme

\oldpage[IV-3]{190}
$B=A[T_1,\ldots,T_n]/\mathfrak J$. Soit $(\mathfrak J_\lambda)$ la famille des
idéaux de type fini de $A[T_1,\ldots,T_n]$ contenus dans $\mathfrak J$, de
sorte que $\mathfrak J$ est la réunion filtrante des $\mathfrak J_\lambda$~;
si $X$ et les
$X_\lambda=\operatorname{Spec}(A[T_1,\ldots,T_n]/\mathfrak J_\lambda)$ sont
considérés comme des sous-préschémas fermés de
$Z=\operatorname{Spec}(A[T_1,\ldots,T_n])$, on a donc, pour les espaces
sous-jacents, $X=\bigcap_\lambda X_\lambda$. Si $f_\lambda:X_\lambda\to S$ est
le morphisme structural, on en déduit que
$f^{-1}(s)=\bigcap_\lambda f_\lambda^{-1}(s)$ pour tout $s\in S$, et comme les
ensembles $f_\lambda^{-1}(s)$ sont fermés dans l'espace noethérien $Z_s$, il
existe un $\lambda$ (dépendant de $s$) tel que
$f^{-1}(s)=f_\lambda^{-1}(s)$. Si alors, pour tout $x\in X$, on pose
$d(x)=\dim_x(f^{-1}(f(x)))$, $d_\lambda(x)=\dim_x(f_\lambda^{-1}(f_\lambda(x)))$,
ce qui précède prouve que l'on a $d(x)=\inf_\lambda d_\lambda(x)$~; les
fonctions $d_\lambda$ étant semi-continues supérieurement d'après la première
partie de la démonstration, il en est de même de $d$. C.Q.F.D.
\end{proof}

\begin{corollary}[13.1.4]
\label{IV.13.1.4-fr}
Sous les hypothèses de (13.1.3), l'ensemble des $x\in X$ tels que $x$ soit
isolé dans $f^{-1}(f(x))$ est ouvert dans $X$.
\end{corollary}

\begin{proof}
En effet, c'est le complémentaire de $F_1(X)$ ($\mathbf 0$, 14.1.10).
\end{proof}

On notera qu'on retrouve ainsi, sous des hypothèses plus générales, la
conséquence (\textbf{III}, 4.4.10) du «~Main theorem~» de Zariski.

\begin{corollary}[13.1.5]
\label{IV.13.1.5-fr}
Sous les hypothèses de (13.1.3), supposons en outre que $f$ soit un morphisme
fermé. Alors, pour tout entier $n$, l'ensemble des $y\in Y$ tels que
$\dim(f^{-1}(y))\geq n$ est fermé, autrement dit, l'application
$y\mapsto\dim(f^{-1}(y))$ est semi-continue supérieurement~; en particulier,
si $y$ est spécialisation de $y'$, on a
$\dim(f^{-1}(y))\geq\dim(f^{-1}(y'))$.
\end{corollary}

\begin{proof}
En effet, dire que $\dim(f^{-1}(y))\geq n$ signifie que $y\in f(F_n(X))$
($\mathbf 0$, 14.1.6).
\end{proof}

\begin{corollary}[13.1.6]
\label{IV.13.1.6-fr}
Soient $X$, $Y$ deux préschémas irréductibles, $f:X\to Y$ un morphisme
dominant localement de type fini. Soit $\xi$ (resp. $\eta$) le point générique
de $X$ (resp. $Y$) et soit $e=\dim(f^{-1}(\eta))$. Alors, pour tout $x\in X$,
on a $\dim_x(f^{-1}(f(x)))\geq e$.
\end{corollary}

\begin{proof}
En effet, l'ensemble des $x$ tels que $\dim_x(f^{-1}(f(x)))<e$ est ouvert en
vertu de (13.1.3), et comme il ne peut contenir $\xi$, il est vide.
\end{proof}

Notons enfin le résultat plus facile suivant~:

\begin{proposition}[13.1.7]
\label{IV.13.1.7-fr}
Soient $Y$ un préschéma quasi-compact, $f:X\to Y$ un morphisme de type fini.
Il existe un entier $n$ tel que, pour tout $y\in Y$, on ait
$\dim(f^{-1}(y))\leq n$.
\end{proposition}

\begin{proof}
Comme il y a un recouvrement ouvert affine fini $(V_i)$ de $Y$ tel que chaque
$f^{-1}(V_i)$ soit réunion finie d'ouverts affines, on est aussitôt ramené au
cas où $Y=\operatorname{Spec}(A)$ et $X=\operatorname{Spec}(B)$ sont affines,
$B$ étant une $A$-algèbre de type fini. Si $B$ admet un système de $n$
générateurs, alors, pour tout $y\in Y$, $B\otimes_Ak(y)$ est une
$k(y)$-algèbre admettant $n$ générateurs, donc $\dim(f^{-1}(y))\leq n$ en
vertu de (4.1.1).
\end{proof}

\subsection{Morphismes équidimensionnels : cas des morphismes dominants de préschémas irréductibles}
\label{subsection:IV.13.2-fr}

\begin{env}[13.2.1]
\label{IV.13.2.1-fr}
Soient $Y$ un préschéma irréductible, $X$ un préschéma irréductible,
$f:X\to Y$ un morphisme dominant localement de type fini~; soit $\eta$ le
point générique de $Y$. On sait (13.1.6) que pour tout $x\in X$, on a
$\dim_x(f^{-1}(f(x)))\geq\dim(f^{-1}(\eta))$.
\end{env}

\oldpage[IV-3]{191}
\begin{definition}[13.2.2]
\label{IV.13.2.2-fr}
Sous les hypothèses de (13.2.1), on dit que $f$ est équidimensionnel au point
$x$ (ou que $X$ est équidimensionnel sur $Y$ au point $x$) si
\[
  \dim_x(f^{-1}(f(x)))=\dim(f^{-1}(\eta)).
\]
On dit que $f$ est équidimensionnel (ou que $X$ est équidimensionnel sur $Y$)
si $f$ est équidimensionnel en tout point $x\in X$.
\end{definition}

Il résulte du théorème de Chevalley (13.1.3) que l'ensemble des $x\in X$ où
$f$ est équidimensionnel est \emph{ouvert non vide}. En outre, si $f$ est
équidimensionnel au point $x$, toutes les composantes irréductibles de
$f^{-1}(f(x))$ qui contiennent $x$ ont \emph{même dimension}, car chacune
d'elles a une dimension qui est $\geq\dim(f^{-1}(\eta))$ (13.1.6) et
$\leq\dim_x(f^{-1}(f(x)))$ en vertu de ($\mathbf 0$, 14.1.5).

\begin{proposition}[13.2.3]
\label{IV.13.2.3-fr}
Soient $Y$ un préschéma irréductible localement noethérien, $X$ un préschéma
irréductible, $f:X\to Y$ un morphisme dominant localement de type fini~;
soient $\eta$ le point générique de $Y$, $x$ un point de $X$, $y=f(x)$, et
supposons que l'on ait
\[
\tag{13.2.3.1}
\label{IV.13.2.3.1-fr}
  \dim(\mathcal O_x)=\dim(\mathcal O_y)
  +\dim(\mathcal O_x\otimes_{\mathcal O_y}k(y)).
\]
Alors $f$ est équidimensionnel au point $x$. La réciproque est vraie si les
deux membres de l'inégalité (5.6.5.2) sont égaux, en particulier si
$\mathcal O_y$ est universellement caténaire.
\end{proposition}

\begin{proof}
Cela résulte aussitôt de (5.6.5.2) et de l'inégalité $\delta(x)\geq0$
(13.1.1).
\end{proof}

\begin{env}[13.2.4]
\label{IV.13.2.4-fr}
Soient maintenant, de façon générale, $Y$ un préschéma localement noethérien,
$f:X\to Y$ un morphisme de type fini, $x$ un point de $X$, $X'$ une partie
fermée irréductible de $X$ contenant $x$, $y=f(x)$,
$Y'=\overline{f(X')}$, $\eta'$ le point générique de $Y'$. Désignons encore
par $X'$ et $Y'$ les sous-préschémas fermés réduits de $X$ et $Y$
respectivement ayant $X'$ et $Y'$ pour espaces sous-jacents~; la restriction
$X'\to Y$ de $f$ se factorise donc en
$X'\xrightarrow{f'}Y'\xrightarrow{j}Y$, où $j$ est l'injection canonique
(\textbf{I}, 5.2.2), et $f'$ est de type fini (1.5.4). Posons alors pour
abréger
\[
\tag{13.2.4.1}
\label{IV.13.2.4.1-fr}
  A=\mathcal O_{Y,y},\quad B=\mathcal O_{X,x},\quad
  A'=\mathcal O_{Y',y},\quad B'=\mathcal O_{X',x}.
\]
La formule (5.6.5.2) appliquée au morphisme dominant $f'$ et aux préschémas
irréductibles $X'$, $Y'$, donne
\[
\tag{13.2.4.2}
\label{IV.13.2.4.2-fr}
  \dim(B')\leq\dim(A')+\dim(B'\otimes_{A'}k(y))
  -(\dim_x(f'^{-1}(y))-\dim(f'^{-1}(\eta'))).
\]
D'autre part, l'anneau local $A'$ (resp. $B'$, $B'\otimes_{A'}k(y)$) est un
anneau quotient de $A$ (resp. $B$, $B\otimes_Ak(y)$), donc
($\mathbf 0$, 16.1.2.1) on a
\[
\tag{13.2.4.3}
\label{IV.13.2.4.3-fr}
  \dim(A')\leq\dim(A),\quad \dim(B')\leq\dim(B),\quad
  \dim(B'\otimes_{A'}k(y))\leq\dim(B\otimes_Ak(y)).
\]
On déduit donc en premier lieu de (13.2.4.3) et ($\mathbf 0$, 16.3.9)
\[
\tag{13.2.4.4}
\label{IV.13.2.4.4-fr}
  \dim(B')\leq\dim(A')+\dim(B'\otimes_{A'}k(y))
  \leq\dim(A)+\dim(B\otimes_Ak(y)).
\]
En outre, en vertu de ($\mathbf 0$, 16.3.9), on a aussi les inégalités
\[
\tag{13.2.4.5}
\label{IV.13.2.4.5-fr}
  \dim(B')\leq\dim(B)\leq\dim(A)+\dim(B\otimes_Ak(y)).
\]
\end{env}

\oldpage[IV-3]{192}
La comparaison de ces inégalités montre donc que~:

\begin{lemma}[13.2.5]
\label{IV.13.2.5-fr}
Avec les notations de (13.2.4), les conditions suivantes sont équivalentes~:
\begin{enumerate}[label=\emph{\alph*)}]
  \item $\dim(B')=\dim(A)+\dim(B\otimes_Ak(y))$.
  \item On a simultanément les relations suivantes~:
  \begin{enumerate}[label=(\roman*)]
    \item $\dim(A')=\dim(A)$.
    \item $\dim(B'\otimes_{A'}k(y))=\dim(B\otimes_Ak(y))$, autrement dit
    \[
      \dim_x(X'\cap f^{-1}(y))=\dim_x(f^{-1}(y)).
    \]
    \item $\dim_x(f'^{-1}(y))=\dim(f'^{-1}(\eta'))$, autrement dit
    $f':X'\to Y'$ est équidimensionnel (13.2.2) au point $x$.
    \item On a l'égalité
    \[
      \dim(B')=\dim(A')+\dim(B'\otimes_{A'}k(y))
      -(\dim_x(f'^{-1}(y))-\dim(f'^{-1}(\eta')))
    \]
  \end{enumerate}
  (relation qui est toujours vérifiée lorsque $A=\mathcal O_{Y,y}$ est un
  anneau universellement caténaire, en vertu de (5.6.5)).
  \item On a simultanément les relations suivantes~:
  \begin{enumerate}[label=(\roman*)]
    \item $\dim(B')=\dim(B)$.
    \item $\dim(B)=\dim(A)+\dim(B\otimes_Ak(y))$.
  \end{enumerate}
\end{enumerate}
\end{lemma}

\begin{env}[13.2.6]
\label{IV.13.2.6-fr}
Rappelons maintenant que les composantes irréductibles $X_i$ de $X$ contenant
$x$ sont en nombre fini et que l'on a ((5.1.2.1) et ($\mathbf 0$, 14.2.1.1))
\[
\tag{13.2.6.1}
\label{IV.13.2.6.1-fr}
  \dim(\mathcal O_{X,x})=\sup_i\dim(\mathcal O_{X_i,x}).
\]
L'équivalence des conditions $b)$ et $c)$ dans (13.2.5) implique par suite,
compte tenu de ($\mathbf 0$, 16.3.9) et ($\mathbf 0$, 14.2.1)~:
\end{env}

\begin{proposition}[13.2.7]
\label{IV.13.2.7-fr}
Soient $Y$ un préschéma localement noethérien, $f:X\to Y$ un morphisme de type
fini, $x$ un point de $X$, $f(x)=y$~; on a
\[
\tag{13.2.7.1}
\label{IV.13.2.7.1-fr}
  \dim(\mathcal O_{X,x})\leq\dim(\mathcal O_{Y,y})
  +\dim(\mathcal O_{X,x}\otimes_{\mathcal O_{Y,y}}k(y)).
\]
Pour que les deux membres de (13.2.7.1) soient égaux, il faut et il suffit
qu'il existe une partie fermée irréductible $X'$ de $X$ contenant $x$ et
vérifiant simultanément les conditions suivantes~:
\begin{enumerate}[label=(\roman*)]
  \item Si $Y'=\overline{f(X')}$, on a
  $\dim(\mathcal O_{Y',y})=\dim(\mathcal O_{Y,y})$.
  \item $\dim_x(X'\cap f^{-1}(y))=\dim_x(f^{-1}(y))$ (autrement dit, $X'$
  contient une des composantes irréductibles de $f^{-1}(y)$ contenant $x$, de
  dimension maxima parmi toutes ces composantes). Il revient au même de dire
  que
  $\dim(\mathcal O_{X',x}\otimes_{\mathcal O_{Y',y}}k(y))
  =\dim(\mathcal O_{X,x}\otimes_{\mathcal O_{Y,y}}k(y))$.
  \item $X'$ est équidimensionnel sur $Y'$ au point $x$, autrement dit, on a
  \[
    \dim_x(X'\cap f^{-1}(y))=\dim(X'\cap f^{-1}(\eta')),
  \]
  où $\eta'$ est le point générique de $Y'$ (et par suite, toutes les
  composantes irréductibles de $X'\cap f^{-1}(y)$ contenant $x$ ont même
  dimension (13.2.2)).

  \oldpage[IV-3]{193}
  \item On a l'égalité
  \[
    \dim(\mathcal O_{X',x})=\dim(\mathcal O_{Y',y})
    +\dim(\mathcal O_{X',x}\otimes_{\mathcal O_{Y',y}}k(y))
  \]
  (condition toujours impliquée par (iii) lorsque $\mathcal O_{Y,y}$ est un
  anneau universellement caténaire).
\end{enumerate}
De plus, $X'$ est alors une composante irréductible de $X$ et $Y'$ une
composante irréductible de $Y$.
\end{proposition}

Dans l'énoncé, les anneaux locaux $\mathcal O_{X',x}$ et
$\mathcal O_{Y',y}$ sont relatifs aux sous-préschémas fermés réduits de $X$,
$Y$ respectivement ayant $X'$ et $Y'$ pour espaces sous-jacents respectifs.

En outre, cela prouve que

\begin{corollary}[13.2.8]
\label{IV.13.2.8-fr}
Si les deux membres de (13.2.7.1) sont égaux, les parties fermées irréductibles
$X'$ de $X$ contenant $x$ et vérifiant les conditions (i) à (iv) de (13.2.7)
sont exactement celles pour lesquelles on a
$\dim(\mathcal O_{X',x})=\dim(\mathcal O_{X,x})$ (ce qui implique
nécessairement que $X'$ est une composante irréductible de $X$).
\end{corollary}

La proposition (13.2.7) entraîne~:

\begin{corollary}[13.2.9]
\label{IV.13.2.9-fr}
Soient $Y$ un préschéma localement noethérien, $f:X\to Y$ un morphisme de type
fini, $x$ un point de $X$, $f(x)=y$. Supposons que $\mathcal O_{Y,y}$ soit un
anneau universellement caténaire. Alors les conditions suivantes sont
équivalentes~:
\begin{enumerate}[label=\emph{\alph*)}]
  \item Les deux membres de (13.2.7.1) sont égaux.
  \item Il existe une composante irréductible $Z$ de $f^{-1}(y)$ contenant
  $x$, de dimension $\dim_x(f^{-1}(y))$, telle que, pour tout $x'$ dans un
  voisinage de $x$ dans $Z$, on ait
  \[
  \tag{13.2.9.1}
  \label{IV.13.2.9.1-fr}
    \dim(\mathcal O_{X,x'})=\dim(\mathcal O_{Y,y})
    +\dim(\mathcal O_{X,x'}\otimes_{\mathcal O_{Y,y}}k(y)).
  \]
  \item Il existe une composante irréductible $Z$ de $f^{-1}(y)$ contenant
  $x$, de dimension $\dim_x(f^{-1}(y))$, telle que pour le point générique $z$
  de $Z$, on ait
  \[
  \tag{13.2.9.2}
  \label{IV.13.2.9.2-fr}
    \dim(\mathcal O_{X,z})=\dim(\mathcal O_{Y,y}).
  \]
\end{enumerate}
\end{corollary}

Montrons que $a)$ entraîne $b)$. Soit $\dim_x(f^{-1}(y))=n$~; en vertu de
(13.2.7), il existe une composante irréductible $X'$ de $X$ vérifiant les
conditions (i) à (iv) de (13.2.7)~; soit $Z$ une composante irréductible de
dimension $n$ de $f^{-1}(y)$, contenant $x$ et contenue dans $X'$. Comme
$f^{-1}(y)$ est localement noethérien, il existe un voisinage ouvert $U$ de
$x$ dans $Z$ tel que $U$ ne rencontre aucune composante irréductible de
$f^{-1}(y)$ autre que celles qui contiennent $x$, donc (4.1.1.3)
$\dim_{x'}(f^{-1}(y))=n$ pour tout $x'\in U$~; il est clair alors que les
conditions (i) à (iii) de (13.2.7) sont satisfaites quand on y remplace $x$
par un point quelconque $x'\in U$, et il en est de même de la condition (iv)
puisque $\mathcal O_{Y,y}$ est universellement caténaire~; d'où la conclusion
par (13.2.7). La condition $b)$ entraîne trivialement $c)$ en vertu de
(5.1.2). Enfin, si $c)$ est vérifiée et si $X''$ est une composante
irréductible de $X$ contenant $Z$ et telle que les conditions (i) à (iv) de
(13.2.7) soient vérifiées quand on y remplace $X'$ par $X''$ et $x$ par $z$,
il est clair que ces conditions sont aussi vérifiées pour $X''$ et $x$ puisque
$\mathcal O_{Y,y}$ est universellement caténaire, donc $c)$ implique $a)$.

\begin{proposition}[13.2.10]
\label{IV.13.2.10-fr}
Soient $Y$ un préschéma irréductible localement noethérien, $\eta$ son point
générique, $f:X\to Y$ un morphisme de type fini, $y$ un point de $f(X)$.
Soient $Z_i$ les

\oldpage[IV-3]{194}
composantes irréductibles de $f^{-1}(y)$, $z_i$ le point générique de $Z_i$,
et considérons les conditions suivantes~:
\begin{enumerate}[label=\emph{\alph*)}]
  \item Pour tout $x\in f^{-1}(y)$, on a la relation
  \[
  \tag{13.2.10.1}
  \label{IV.13.2.10.1-fr}
    \dim(\mathcal O_{X,x})=\dim(\mathcal O_{Y,y})
    +\dim(\mathcal O_{X,x}\otimes_{\mathcal O_{Y,y}}k(y)).
  \]
  \item Pour tout $i$, on a
  \[
  \tag{13.2.10.2}
  \label{IV.13.2.10.2-fr}
    \dim(\mathcal O_{X,z_i})=\dim(\mathcal O_{Y,y}).
  \]
  \item Pour tout $i$, il existe une composante irréductible $X_i$ de $X$
  contenant $Z_i$ et telle que
  $\dim(X_i\cap f^{-1}(\eta))=\dim(Z_i)$ (autrement dit, telle que le
  \emph{sous-préschéma fermé réduit} $X_i$ de $X$ soit équidimensionnel sur $Y$
  au point $z_i$).
\end{enumerate}
Alors $a)$ entraîne $b)$ et $b)$ entraîne $c)$~; en outre, si
$\mathcal O_{Y,y}$ est universellement caténaire, les trois conditions $a)$,
$b)$, $c)$ sont équivalentes.
\end{proposition}

L'anneau $\mathcal O_{X,z_i}\otimes_{\mathcal O_{Y,y}}k(y)$ étant de dimension
$0$ (5.1.2), $a)$ entraîne évidemment $b)$~; $b)$ entraîne $c)$ en vertu de
(13.2.7) appliqué au point $z_i$. Inversement, supposons que
$\mathcal O_{Y,y}$ soit universellement caténaire~; comme la condition $c)$
implique que $f(X_i)$ est dense dans $Y$, il résulte de $c)$ que les conditions
(i) à (iv) de (13.2.7) sont vérifiées en remplaçant $X'$ par $X_i$ et $x$ par
$z_i$, donc $c)$ implique $b)$~; enfin $b)$ implique $a)$ en vertu de
(13.2.9).

\begin{corollary}[13.2.11]
\label{IV.13.2.11-fr}
Les notations étant celles de (13.2.10), supposons que $\mathcal O_{Y,y}$ soit
universellement caténaire. Pour tout $y'\in Y$, soit $E(y')$ l'ensemble des
dimensions des composantes irréductibles de $f^{-1}(y')$ et posons
$d(y')=\dim(f^{-1}(y'))=\sup(E(y'))$. Alors, si les conditions équivalentes
$a)$, $b)$, $c)$ de (13.2.10) sont vérifiées, on a
$E(y)\subset E(\eta)$, d'où $d(y)\leq d(\eta)$.
\end{corollary}

En effet, avec les notations de (13.2.10), $X_i\cap f^{-1}(\eta)$ n'est pas
vide et est par suite une composante irréductible de $f^{-1}(\eta)$
($\mathbf 0_{\mathrm I}$, 2.1.8).

\begin{remarks}[13.2.12]
\label{IV.13.2.12-fr}
\begin{enumerate}[label=(\roman*)]
  \item Rappelons (6.1.2) que la relation (13.2.10.1) est toujours vérifiée
  lorsque le morphisme $f$ est \emph{plat} au point $x$.
  \item Avec les notations de (13.2.10), supposons que $\mathcal O_{Y,y}$ soit
  universellement caténaire et en outre que le morphisme $f$ soit \emph{propre}~;
  alors, si les conditions équivalentes $a)$, $b)$, $c)$ de (13.2.10) sont
  satisfaites, on a même $d(y)=d(\eta)$, car il résulte de (13.1.5) que l'on a
  $d(y)\geq d(\eta)$.
  \item Le morphisme de (12.2.3, (ii)) est \emph{propre} et \emph{plat} et tous
  les anneaux locaux de $Y$ sont universellement caténaires~; en outre, les deux
  composantes irréductibles $X_1$, $X_2$ de $X$ sont \emph{équidimensionnelles}
  sur $Y$ en tout point~; mais $E(y)$ a \emph{deux} éléments pour tout
  $y\neq y_0$, alors que $E(y_0)$ est réduit à un \emph{seul} élément, donc
  $E(y)$ n'est \emph{pas constant} dans $Y$.
\end{enumerate}
\end{remarks}

\subsection{Morphismes équidimensionnels : cas général}
\label{subsection:IV.13.3-fr}

\begin{proposition}[13.3.1]
\label{IV.13.3.1-fr}
Soient $Y$ un préschéma, $f:X\to Y$ un morphisme localement de type fini,
$x$ un point de $X$, $y=f(x)$. Désignons par $Y_\alpha$ les composantes
irréductibles de $Y$ contenant $y$. Alors les conditions suivantes sont
équivalentes~:
\begin{enumerate}[label=\emph{\alph*)}]
\oldpage[IV-3]{195}
  \item Il existe un entier $e$ et un voisinage ouvert $U$ de $x$ tels que
  l'image par $f$ de toute composante irréductible de $U$ soit dense dans un
  $Y_\alpha$, et que, pour tout $x'\in U$, l'espace
  $U\cap f^{-1}(f(x'))$ soit équidimensionnel et de dimension $e$.
  \item[\emph{a$'$)}] Il existe un entier $e$ et un voisinage ouvert $U$ de
  $x$ tels que l'image par $f$ de toute composante irréductible de $U$ soit
  dense dans un $Y_\alpha$ et que, si on désigne par $y_\alpha$ le point
  générique de $Y_\alpha$, toute composante irréductible des espaces
  $U\cap f^{-1}(y)$, $U\cap f^{-1}(y_\alpha)$ soit de dimension $e$.
  \item[\emph{a$''$)}] Il existe un entier $e$ et un voisinage ouvert $U$ de
  $x$ tels que, pour chacune des composantes irréductibles $U_\lambda$ de
  $U$, $f(U_\lambda)$ soit dense dans un $Y_\alpha$ et que, pour tout
  $x'\in U_\lambda$, les composantes irréductibles de
  $U_\lambda\cap f^{-1}(f(x'))$ soient toutes de dimension $e$.
  \item Il existe un entier $e$, un voisinage ouvert $U$ de $x$ et un
  $Y$-morphisme quasi-fini
  \[
    g:U\longrightarrow Y\otimes_{\mathbf Z}\mathbf Z[T_1,\ldots,T_e]
  \]
  (préschéma que nous noterons aussi $Y[T_1,\ldots,T_e]$ pour abréger) tels
  que l'image par $g$ de toute composante irréductible de $U$ soit dense dans
  une composante irréductible de $Y[T_1,\ldots,T_e]$.
\end{enumerate}
\end{proposition}

\begin{proof}
Il est immédiat que $a''$) entraîne $a$), car pour tout $x\in U$, les
composantes irréductibles de $U\cap f^{-1}(f(x))$ sont chacune composante
irréductible d'un des $U_\lambda\cap f^{-1}(f(x))$, d'où la conclusion par
($\mathbf 0$, 14.1.4). La condition $a$) entraîne trivialement $a'$).
Montrons ensuite que $a'$) entraîne $a''$)~; on peut se borner au cas où $f$
est de type fini et $X$ et $Y$ réduits (1.5.4). Soit $U_\lambda$ une
composante irréductible de $U$, et supposons que $f(U_\lambda)$ soit dense
dans $Y_\alpha$~; alors la restriction de $f$ à $U_\lambda$ se factorise en
$U_\lambda\xrightarrow{f_\alpha}Y_\alpha\to Y$, où $f_\alpha$ est de type
fini et dominant (\textbf{I}, 5.2.2). Soit $x_\lambda$ le point générique de
$U_\lambda$~; en vertu de ($\mathbf 0_{\mathrm I}$, 2.1.8),
$U_\lambda\cap f_\alpha^{-1}(y_\alpha)$ est l'unique composante irréductible
de $U\cap f^{-1}(y_\alpha)$ contenant $x_\lambda$ et est par hypothèse de
dimension $e$, égale aux dimensions de toutes les composantes irréductibles
de $U_\lambda\cap f^{-1}(y)$, en vertu de l'hypothèse et de (13.1.1). Mais en
vertu du théorème de Chevalley (13.1.3) et de (13.1.1), l'ensemble $V_\lambda$
des $x'\in U_\lambda$ tels que
$\dim_{x'}(f_\lambda^{-1}(f_\lambda(x')))=e$ est ouvert et contient $x$, et il
suffit de prendre la réunion des $V_\lambda$ pour obtenir un ensemble ouvert
satisfaisant aux conditions de $a''$).

Prouvons maintenant que $a$) entraîne $b$)~; on peut se limiter au cas où
$Y=\operatorname{Spec}(A)$ et $X=\operatorname{Spec}(B)$ sont affines et où
$U=X$. Prouvons d'abord le lemme suivant~:

\begin{lemma}[13.3.1.1]
\label{IV.13.3.1.1-fr}
Soient $Y=\operatorname{Spec}(A)$, $X=\operatorname{Spec}(B)$ deux schémas
affines, $f:X\to Y$ un morphisme de type fini, $y$ un point de $f(X)$, et
posons $e=\dim(f^{-1}(y))$. Alors il existe un $Y$-morphisme
$g:X\to Y[T_1,\ldots,T_e]$ tel que (si l'on pose
$X_y=X\times_Y\operatorname{Spec}(k(y))$), le morphisme
$g_y:X_y\to\operatorname{Spec}(k(y)[T_1,\ldots,T_e])$ soit fini. En outre,
pour un tel morphisme $g$, $g_y$ est nécessairement surjectif et il existe un
voisinage ouvert $U$ de $f^{-1}(y)$ dans $X$ tel que $g|U$ soit quasi-fini.
\end{lemma}

Soit $\mathfrak p=\mathfrak j_y$~; l'anneau $B\otimes_Ak(\mathfrak p)$ est une
$k(\mathfrak p)$-algèbre de type fini, donc le lemme de normalisation
(Bourbaki, \emph{Alg. comm.}, chap.~V, \S~3, n\textsuperscript{o}~1, th.~1)
prouve qu'il y a dans $B\otimes_Ak(\mathfrak p)$ une suite finie
$(t_i)_{1\leq i\leq r}$ d'éléments algébriquement indépendants sur
$k(\mathfrak p)$ et tels que si l'on pose
$C'=k(\mathfrak p)[t_1,\ldots,t_r]$, $B\otimes_Ak(\mathfrak p)$ soit une
$C'$-algèbre \emph{finie}~; on a donc
$\dim(B\otimes_Ak(\mathfrak p))=\dim(C')$ ($\mathbf 0$, 16.1.5) et comme
$\dim(C')=r$ (5.2.1), on a $r=e$. Comme
$B\otimes_Ak(\mathfrak p)=(B/\mathfrak pB)\otimes_{A/\mathfrak p}k(\mathfrak p)$,
on peut, en multipliant les $t_i$ par un élément $\neq0$

\oldpage[IV-3]{196}
convenable de $A/\mathfrak p$, supposer que chaque $t_i$ est l'image canonique
dans $B\otimes_Ak(\mathfrak p)$ d'un élément $s_i\in B$. Soit alors
$u:A[T_1,\ldots,T_e]\to B$ l'homomorphisme tel que $u(T_i)=s_i$ pour tout
$i$, et soit $g:X\to Y[T_1,\ldots,T_e]$ le morphisme correspondant. Il est
clair qu'en raison du choix des $s_i$, $g_y$ est un morphisme fini. Pour tout
morphisme $g:X\to Y[T_1,\ldots,T_e]$, tel que $g_y$ soit fini, il résulte de
(5.4.2) et (4.1.2.1) que $g_y$ est nécessairement surjectif. D'autre part, en
vertu de (13.1.4), l'ensemble $U$ des $x\in X$ qui sont isolés dans leur fibre
$g^{-1}(g(x))$ est ouvert dans $X$ et contient $f^{-1}(y)$, et par définition
la restriction $g|U$ est un morphisme quasi-fini.

Ce lemme étant établi, pour prouver que $a$) implique $b$), il reste à voir
que si $x_\lambda$ est un point maximal de $U$, $g(x_\lambda)=z_\lambda$ est
un point maximal de $Z=Y[T_1,\ldots,T_e]$. En vertu de l'hypothèse $a$), on
peut (en restreignant au besoin $U$), supposer que $f(x_\lambda)$ est l'un des
points génériques $y_\alpha$ des composantes irréductibles de $Y$ contenant
$y$~; si $p:Z\to Y$ est le morphisme structural, on a donc
$p(z_\lambda)=y_\alpha$, et l'on déduit par suite de $g$ un
$k(y_\alpha)$-morphisme quasi-fini
\[
  h:U\cap f^{-1}(y_\alpha)\longrightarrow p^{-1}(y_\alpha).
\]
Or $p^{-1}(y_\alpha)=\operatorname{Spec}(k(y_\alpha)[T_1,\ldots,T_e])$ est
intègre et de dimension $e$~; si $z_\lambda$ n'était pas point maximal de
$Z$, il ne serait pas point maximal de $p^{-1}(y_\alpha)$
($\mathbf 0_{\mathrm I}$, 2.1.8) et son adhérence $Z'_\lambda$ dans
$p^{-1}(y_\alpha)$ serait donc de dimension $<e$ (4.1.2.1). Mais comme $h$
est quasi-fini, il résulte de (4.1.2) et de l'hypothèse $a$) que l'on a
$\dim(Z'_\lambda)\geq e$ (la restriction de $h$ à
$\overline{\{x_\lambda\}}$ se factorisant en
\[
  \overline{\{x_\lambda\}}\longrightarrow Z'_\lambda
  \longrightarrow p^{-1}(y_\alpha)
\]
en vertu de (\textbf{I}, 5.2.2))~; on aboutit ainsi à une contradiction, ce
qui montre que $a$) entraîne $b$).

Prouvons finalement que $b$) implique $a$). Notons que le morphisme structural
$p:Z=Y[T_1,\ldots,T_e]\to Y$ est fidèlement plat~; donc (2.3.4) les points
maximaux de $Z$ ont pour images par $p$ les points maximaux de $Y$~; ceci
prouve déjà que les $f(x_\lambda)$ sont points génériques des $Y_\alpha$. En
outre, si $f(x_\lambda)=y_\alpha$, le $k(y_\alpha)$-morphisme
$h:U\cap f^{-1}(y_\alpha)\to p^{-1}(y_\alpha)$ déduit de $g$ est dominant et
quasi-fini par hypothèse~; on a donc
$\dim(U_\lambda\cap f^{-1}(y_\alpha))=e$ en vertu de (4.1.2)
($U_\lambda$ étant la composante irréductible de $U$ de point générique
$x_\lambda$). De même, pour tout $x'\in U_\lambda$, le morphisme
$U_\lambda\cap f^{-1}(f(x'))\to p^{-1}(f(x'))$ déduit de $g$ est un
$k(f(x'))$-morphisme quasi-fini, donc (4.1.2) on a
$\dim(f^{-1}(f(x'))\cap U_\lambda)\leq e$~; mais par ailleurs on sait (13.1.6)
que toutes les composantes irréductibles de
$U_\lambda\cap f^{-1}(f(x'))$ sont de dimension $\geq e$~; on voit ainsi que
ces composantes sont exactement de dimension $e$, donc $b$) entraîne $a$).
C.Q.F.D.
\end{proof}

\begin{definition}[13.3.2]
\label{IV.13.3.2-fr}
Soient $Y$ un préschéma, $f:X\to Y$ un morphisme localement de type fini, $x$
un point de $X$. On dit que $f$ est équidimensionnel au point $x$ (ou que $X$
est équidimensionnel sur $Y$ au point $x$) si les conditions équivalentes de
(13.3.1) sont vérifiées. On dit que $f$ est équidimensionnel (ou que $X$ est
équidimensionnel sur $Y$) si $f$ est équidimensionnel en tout point $x\in X$.
\end{definition}

Il est clair que, pour que $f$ soit équidimensionnel en un point $x$, il faut
et il suffit que $f_{\mathrm{red}}$ le soit. En outre, les conditions de
(13.3.1) montrent que l'ensemble des points où $f$ est équidimensionnel est
\emph{ouvert} dans $X$.

On notera que lorsque $X$ et $Y$ sont \emph{irréductibles}, dire que $f$ est
équidimen-
\oldpage[IV-3]{197}
sionnel au point $x$ signifie, en vertu de (13.3.1, $a''$)), que $f$ est
dominant et que
$\dim_x(f^{-1}(f(x)))=\dim(f^{-1}(\eta))$ (où $\eta$ est le point générique de
$Y$)~; la définition (13.3.2) coïncide donc dans ce cas avec la définition
(13.2.2).

\begin{proposition}[13.3.3]
\label{IV.13.3.3-fr}
Soient $Y$ un préschéma localement noethérien, $f:X\to Y$ un morphisme
localement de type fini, $x$ un point de $X$, $X_j$ les composantes
irréductibles de $X$ (en nombre fini) contenant $x$. Pour que $f$ soit
équidimensionnel au point $x$, il faut et il suffit que, pour tout $j$,
$Y_j=\overline{f(X_j)}$ soit une composante irréductible de $Y$ et, en
désignant par $X_j$ et $Y_j$ les sous-préschémas fermés réduits de $X$ et $Y$
d'espaces sous-jacents $X_j$ et $Y_j$, par $f_j:X_j\to Y_j$ le morphisme
déduit de $f$ (\textbf{I}, 5.2.2), par $y_j$ le point générique de $Y_j$, que
$f_j$ soit équidimensionnel au point $x$ et que tous les nombres
$\dim(f_j^{-1}(y_j))$ soient égaux.
\end{proposition}

Cela résulte aussitôt de (13.3.1, $a'$)).

\begin{corollary}[13.3.4]
\label{IV.13.3.4-fr}
Avec les notations de (13.3.3), posons $y=f(x)$. Si $\mathcal O_{Y,y}$ est un
anneau universellement caténaire et si $f$ est équidimensionnel au point $x$,
on a
\[
\tag{13.3.4.1}
\label{IV.13.3.4.1-fr}
  \dim(\mathcal O_{X_j,x})=\dim(\mathcal O_{Y_j,y})+e
  -\deg.\operatorname{tr}_{k(y)}k(x)
\]
où $e$ est la valeur commune des nombres $\dim(f_j^{-1}(y_j))$.
\end{corollary}

Comme chacun des $\mathcal O_{Y_j,y}$, quotient de $\mathcal O_{Y,y}$, est un
anneau universellement caténaire (5.6.1), l'égalité est conséquence de (5.6.5).

\begin{corollary}[13.3.5]
\label{IV.13.3.5-fr}
Avec les notations de (13.3.4), supposons que $f$ soit équidimensionnel au
point $x$ et que $\mathcal O_{Y,y}$ soit un anneau universellement caténaire.
Si l'anneau $\mathcal O_{Y,y}$ est équidimensionnel, il en est de même de
$\mathcal O_{X,x}$~; la réciproque est vraie si l'image par $f$ de la réunion
des $X_j$ est dense dans un voisinage de $y$.
\end{corollary}

En effet, dire que $\mathcal O_{Y,y}$ est équidimensionnel signifie que les
nombres $\dim(\mathcal O_{Y'_i,y})$ sont égaux pour toutes les composantes
irréductibles $Y'_i$ de $Y$ contenant $y$, comme il résulte de (5.1.1.5)
appliqué au schéma local $\operatorname{Spec}(\mathcal O_{Y,y})$~; comme on a
la relation (13.3.4.1), il en résulte par le même raisonnement que
$\mathcal O_{X,x}$ est alors équidimensionnel. Inversement, si
$\mathcal O_{X,x}$ est équidimensionnel, tous les nombres
$\mathcal O_{Y_j,y}$ sont égaux par (13.3.4.1)~; on en déduit que
$\mathcal O_{Y,y}$ est équidimensionnel si les $Y_j$ sont \emph{toutes} les
composantes irréductibles de $Y$ contenant $y$, ce qui résulte de l'hypothèse
supplémentaire.

\begin{proposition}[13.3.6]
\label{IV.13.3.6-fr}
Soient $Y$ un préschéma localement noethérien, $f:X\to Y$ un morphisme
localement de type fini, $x$ un point de $X$, $y=f(x)$. Supposons que l'anneau
$\mathcal O_y$ soit équidimensionnel. Alors, si $\mathcal O_x$ est
équidimensionnel et si l'on a l'égalité
\[
\tag{13.3.6.1}
\label{IV.13.3.6.1-fr}
  \dim(\mathcal O_x)=\dim(\mathcal O_y)
  +\dim(\mathcal O_x\otimes_{\mathcal O_y}k(y))
\]
(cf. (13.2.7.1)) $f$ est équidimensionnel au point $x$, et la réciproque est
vraie si $\mathcal O_y$ est un anneau universellement caténaire.
\end{proposition}

Gardons les notations de (13.3.3)~; il résulte de (13.2.8) et de l'hypothèse
que $\mathcal O_x$ est équidimensionnel, que chacun des $Y_j$ est une
composante irréductible de $Y$ et que chacun des $f_j$ est équidimensionnel au
point $x$~; en outre (13.2.8), on a (compte tenu de (5.6.5))
\[
  \dim(\mathcal O_{X_j,x})=\dim(\mathcal O_{Y_j,y})
  +\dim_x(X_j\cap f^{-1}(y))-\deg.\operatorname{tr}_{k(y)}k(x).
\]

\oldpage[IV-3]{198}
Or, comme $\mathcal O_y$ est supposé équidimensionnel, cette égalité s'écrit
\[
\tag{13.3.6.2}
\label{IV.13.3.6.2-fr}
  \dim_x(X_j\cap f^{-1}(y))=\dim(\mathcal O_{X,x})
  -\dim(\mathcal O_{Y,y})+\deg.\operatorname{tr}_{k(y)}k(x).
\]
Le premier membre de (13.3.6.2) est donc indépendant de $j$~; mais comme $f_j$
est équidimensionnel au point $x$, on a
$\dim_x(X_j\cap f^{-1}(y))=\dim(f_j^{-1}(y_j))$, donc le critère (13.3.3)
montre que $f$ est équidimensionnel au point $x$.

Inversement, supposons que $\mathcal O_{Y,y}$ soit un anneau universellement
caténaire et que $f$ soit équidimensionnel au point $x$~; alors (13.3.5)
$\mathcal O_{X,x}$ est équidimensionnel, et il résulte alors de (13.3.4) que
l'on a la relation
\[
  \dim(\mathcal O_{X,x})=\dim(\mathcal O_{Y,y})+e
  -\deg.\operatorname{tr}_{k(y)}k(x)
\]
où $e$ est la valeur commune des nombres
$\dim(f_j^{-1}(y_j))=\dim_x(X_j\cap f^{-1}(y))$~; mais par définition $e$ est
aussi égal à $\dim_x(f^{-1}(y))$, d'où la relation (13.3.6.1), compte tenu de
(5.6.5.2).

\begin{proposition}[13.3.7]
\label{IV.13.3.7-fr}
Soient $Y$ un préschéma, $f:X\to Y$ un morphisme localement de type fini et
équidimensionnel, $Z$ une partie fermée de $X$. Alors la fonction
\[
\tag{13.3.7.1}
\label{IV.13.3.7.1-fr}
  x\longmapsto\operatorname{codim}_x
  (Z\cap f^{-1}(f(x)),f^{-1}(f(x)))
\]
est semi-continue inférieurement dans $X$.
\end{proposition}

Comme toutes les composantes irréductibles de $f^{-1}(f(x))$ contenant $x$
ont par hypothèse la même dimension (13.3.1), on a, en vertu de (5.2.1),
\[
  \operatorname{codim}_x(Z\cap f^{-1}(f(x)),f^{-1}(f(x)))
  =\dim_x(f^{-1}(f(x)))-\dim_x(Z\cap f^{-1}(f(x))).
\]
Mais par hypothèse le premier terme du second membre est fonction
\emph{continue} de $x$ (13.3.1) et le second est fonction semi-continue
supérieurement de $x$ en vertu de (13.3.3)~; d'où la conclusion.

\begin{proposition}[13.3.8]
\label{IV.13.3.8-fr}
Soient $Y$ un préschéma, $f:X\to Y$ un morphisme localement de type fini, $x$
un point de $X$. Soient $Y'$ un préschéma, $g:Y'\to Y$ un morphisme plat,
$X'=X\times_Y Y'$, $f'=f_{(Y')}:X'\to Y'$~; si $f$ est équidimensionnel au
point $x$, alors $f'$ est équidimensionnel en tout point $x'\in X'$ au-dessus
de $x$.
\end{proposition}

La question étant locale sur $X$ et $Y$, on peut se borner au cas où toute
composante irréductible de $X$ (resp. $Y$) contient $x$ (resp. $y$)~; comme
l'image par $f$ de toute composante irréductible de $X$ est alors dense dans
une composante irréductible de $Y$ (13.3.1), on sait (2.3.5) que l'image par
$f'$ de toute composante irréductible de $X'$ est dense dans une composante
irréductible de $Y'$. D'autre part, par transitivité des fibres
(\textbf{I}, 3.6.4), il résulte de (4.2.8) et de (13.3.1) que, pour tout
$z'\in X'$, l'ensemble des dimensions des composantes irréductibles de
${f'}^{-1}(f'(z'))$ est le même que l'ensemble des dimensions des composantes
irréductibles de $f^{-1}(f(z))$, où $z$ est la projection de $z'$ dans $X$~;
d'où la conclusion en vertu de (13.3.1, $a$)).

\begin{remark}[13.3.9]
\label{IV.13.3.9-fr}
La propriété d'équidimensionnalité d'un morphisme $f:X\to Y$ n'est \emph{pas}
stable par changement de base quelconque $g:Y'\to Y$, même lorsque $g$ est
l'injection canonique d'une composante irréductible $Y'$ de $Y$ dans $Y$
($Y'$ étant

\oldpage[IV-3]{199}
considérée comme sous-préschéma fermé réduit de $Y$). Par exemple, soient $k$
un corps, $A_0=k[S,T]$ l'anneau de polynômes à deux indéterminées,
$\mathfrak a=\mathfrak p_1\mathfrak p_2$, où $\mathfrak p_1$ et
$\mathfrak p_2$ sont les idéaux premiers $A_0S$ et $A_0T$ de $A_0$~; soient
$A=A_0/\mathfrak a$ et $Y=\operatorname{Spec}(A)$, qui a deux composantes
irréductibles $Y_1=\operatorname{Spec}(A/\mathfrak p_1)$,
$Y_2=\operatorname{Spec}(A/\mathfrak p_2)$~; prenons $X=Y_1$,
$f:X\to Y$ étant l'injection canonique, qui est évidemment un morphisme
équidimensionnel. Prenons d'autre part $Y'=Y_2$, $g:Y'\to Y$ étant l'injection
canonique~; alors on a
$X'=X\times_Y Y'=\operatorname{Spec}((A/\mathfrak p_1)\otimes_A
(A/\mathfrak p_2))=\operatorname{Spec}(A/(\mathfrak p_1+\mathfrak p_2))$,
et $\mathfrak p_1+\mathfrak p_2$ est un idéal maximal de $A$, donc $X'$ est
réduit à un point, et le morphisme $f':X'\to Y'$ n'est \emph{pas dominant},
l'image par $f'$ de l'unique point de $X'$ étant un point fermé de $Y'$~;
donc $f'$ n'est pas équidimensionnel.

On peut aussi donner un contre-exemple où $X$ et $Y$ sont \emph{intègres}, $f$
\emph{fini} et \emph{birationnel} (et \emph{a fortiori} équidimensionnel par
(13.3.1, $b$)), $g:Y'\to Y$ \emph{fini et dominant}. Soient $A$ et
$\overline A$ les anneaux locaux définis dans (\textbf{II}, 7.5), et prenons
$Y=\operatorname{Spec}(A)$, $X=\operatorname{Spec}(\overline A)$~; d'autre
part, avec les notations de (\textbf{II}, 7.5), prenons
$Y'=\operatorname{Spec}(\overline B)$~; alors
$X'=\operatorname{Spec}(\overline A\otimes_A\overline B)
=\operatorname{Spec}(\overline B\otimes_B\overline B)$~; mais on vérifie
aussitôt que $\overline B\otimes_B\overline B$ est composé direct des anneaux
$B/\mathfrak p'$, $B/\mathfrak p''$ et de deux anneaux isomorphes à
$B/\mathfrak n$, dont les spectres sont donc réduits à un point~; comme les
projections de ces points sont des points fermés de $Y'$, ici encore on voit
que $f'$ ne transforme pas une composante irréductible de $X'$ en une partie
partout dense d'une composante irréductible de $Y'$, donc $f'$ n'est pas
équidimensionnel.

Dans cet exemple, l'anneau $A$ n'est pas géométriquement unibranche~; nous
allons voir au paragraphe suivant (14.4.6) que de tels phénomènes ne peuvent se
produire lorsque les points de $Y$ sont \emph{géométriquement unibranches}. Le
manque de stabilité de la notion de morphisme équidimensionnel restreint
grandement son intérêt, au profit de la notion de morphisme universellement
ouvert, qui va être étudiée en détail au paragraphe suivant.
\end{remark}

% IV3_B28_P199_END_BEFORE_SECTION_14
